\documentclass[11pt,a4paper]{article}

\usepackage[utf8]{inputenc}

\usepackage[T1]{fontenc}

\usepackage[polish]{babel}

\usepackage{lmodern}

\usepackage{geometry}

\geometry{a4paper,margin=27mm}

\usepackage{microtype}

\usepackage[hidelinks]{hyperref}

% Portfolio: MSC-AAL-P07; source: prospective ARC-XR research scope.

\title{Pomoc przestrzenna AR wobec interfejsu 2D w wykonywaniu bezpiecznych czynności domowych}

% scientific_owner: Mateusz Pomianek (MatPomGit).
% Afiliacja i ORCID autora: pola do uzupelnienia po weryfikacji
% przed ewentualnym zgloszeniem. Bez domyslnych identyfikatorow.
% affiliation_verified:
% orcid_verified:
\author{Mateusz Pomianek}

\date{Wersja 0.6.0, 23 listopada 2026}

\begin{document}

\maketitle

\noindent\textbf{Rodzaj opracowania:} koncepcja przyszłego badania porównawczego (\textit{concept}; status \textit{deferred}).

\begin{abstract}
Rozważono projekt porównania przestrzennych wskazówek rzeczywistości rozszerzonej (AR) z instrukcjami ekranowymi 2D w krótkich sekwencjach neutralnych czynności domowych o minimalnym ryzyku. Celem planowanego badania byłoby ustalenie, czy sposób prezentacji tej samej instrukcji wpływa na poprawność wykonania poszczególnych kroków, zwłaszcza wymagających orientacji przestrzennej. Kandydackim wynikiem głównym jest różnica prawdopodobieństwa poprawnego wykonania kroku między warunkami AR i 2D, z uwzględnieniem zależności obserwacji w obrębie uczestnika i sekwencji. Proponuje się również rejestrowanie pominięć, czasu wykonania i przerwań zadań. Szkic określa hipotezy, podstawowe elementy operacjonalizacji, potencjalne czynniki zakłócające oraz kierunek analizy statystycznej. Projekt ma status \textit{deferred}: nie przeprowadzono w jego ramach udokumentowanego badania, a zatwierdzenie protokołu, przegląd źródeł, weryfikacja wykonalności i rozstrzygnięcia etyczne pozostają otwarte. Nie przedstawia się wyników empirycznych ani dowodów skuteczności lub bezpieczeństwa zastosowań AAL.
\end{abstract}

\noindent\textbf{Słowa kluczowe:} rzeczywistość rozszerzona; instrukcje przestrzenne; interfejs 2D; poprawność wykonywania zadań; czynności życia codziennego; projekt badania porównawczego.

\section{Status i własność naukowa}
MSC-AAL-P07 jest koncepcją badawczą o statusie \textit{deferred} (historycznie \texttt{NEW\_CROSS\_DOMAIN\_IDEA\_HOLD}), a nie ukończonym badaniem. Właścicielem naukowym jest Mateusz Pomianek; kanoniczne źródło tekstu znajduje się w \texttt{MatPomGit/arc-xr}. Własność naukowa, lokalizacja repozytorium i autorstwo są odrębnymi pojęciami. Nie udokumentowano danych empirycznych, prerejestracji ani zgody etycznej.

\section{Wprowadzenie i uzasadnienie}
Instrukcje wykonywania czynności wieloetapowych przekazują co najmniej dwa rodzaje informacji: co należy zrobić i gdzie należy wykonać działanie. W interfejsie ekranowym 2D odbiorca może potrzebować przełożyć opis lub obraz na położenie obiektu w otoczeniu. Przestrzenna wskazówka AR ma natomiast możliwość powiązania komunikatu z miejscem wykonania czynności. Taka różnica formy prezentacji uzasadnia porównanie, ale nie stanowi dowodu wyższej skuteczności AR.

\subsection{Proponowane mechanizmy poznawcze}
Pierwszym kandydackim mechanizmem jest ograniczenie konieczności przenoszenia informacji między ekranem a przestrzenią zadania. Jeśli wskazówka trafnie wskazuje obiekt i jego położenie, może zmniejszać wymagania związane z odnajdywaniem punktu odniesienia. Drugi mechanizm dotyczy uwagi: bezpośrednie odniesienie do obiektu mogłoby ograniczać przełączanie spojrzenia między instrukcją i stanowiskiem. Trzeci wiąże się z zewnętrznym podtrzymywaniem informacji o bieżącym kroku, co hipotetycznie mogłoby ograniczać pominięcia w procedurze wieloetapowej. Są to wyjaśnienia projektowe wymagające osobnej weryfikacji, a nie ustalone mechanizmy działania tego systemu.

\subsection{Mechanizmy konkurencyjne i warunki brzegowe}
AR może również zwiększać obciążenie percepcyjne przez zasłanianie elementów sceny, rozproszenie uwagi, nieprecyzyjne zakotwiczenie wskazówki, opóźnienie jej aktualizacji albo konieczność obsługi dodatkowego urządzenia. Instrukcja 2D może okazać się korzystniejsza w zadaniach czysto sekwencyjnych, gdy położenie przedmiotów nie jest istotne, albo gdy użytkownik już zna stanowisko. Efekt nowości technologii, uczenie się procedury oraz odmienna ergonomia urządzeń mogą prowadzić do pozornej przewagi jednego z warunków. Dlatego dokładność rejestracji przestrzennej, czytelność wskazówek, kolejność prób i uprzednie doświadczenie są potencjalnymi czynnikami zakłócającymi lub modyfikatorami efektu, a nie automatycznie wyjaśnieniami przyczynowymi.

Propozycja dotyczy kontrolowanych czynności laboratoryjnych o minimalnym ryzyku, takich jak manipulowanie nieszkodliwymi przedmiotami. Nie zakłada ekspozycji na urządzenia elektryczne, wysoką temperaturę ani potencjalnie niebezpieczne prace domowe. Trafność ekologiczna i możliwość przeniesienia wyników na wsparcie samodzielności w środowisku domowym pozostają osobnymi pytaniami.

\section{Cel, pytania i zakres nowości}
Celem proponowanej pracy jest porównanie poprawności wykonywania kroków w warunkach informacyjnie równoważnych instrukcji AR i ekranowych instrukcji 2D. Podstawowym przedmiotem wnioskowania ma być wpływ formy wskazówki na zachowanie w określonej klasie bezpiecznych zadań, nie skuteczność kliniczna ani bezpieczeństwo urządzeń.

\paragraph{PB1 (główne).} Czy prawdopodobieństwo poprawnego wykonania kroku różni się między przestrzennymi wskazówkami AR a instrukcjami 2D przy porównywalnej treści i trudności zadań?
\paragraph{PB2 (eksploracyjne).} Czy zależność między warunkiem interfejsu a pominięciami różni się między krokami wymagającymi lokalizacji przestrzennej i krokami wyłącznie sekwencyjnymi?
\paragraph{PB3 (eksploracyjne).} Czy ewentualnym różnicom w poprawności towarzyszą różnice czasu wykonania lub częstości przerywania sekwencji?

Kandydacki wkład naukowy dotyczy mierzalnego kontrastu AR--2D na poziomie poprawności kroku z uwzględnieniem jego charakteru przestrzennego. Samo zastosowanie AR, modeli wizualnych, EEG lub okulometrii nie stanowi dowodu nowości. Ostateczną wartość wkładu należy ocenić po systematyczniejszym przeglądzie badań porównujących te same typy zadań i punkty końcowe.

\section{Stan badań i luka badawcza}
Zhang i współautorzy porównali w zadaniach precyzyjnej manipulacji wskazówki AR o charakterze wizualnym i multisensorycznym, analizując wykonanie i doświadczenie użytkownika \cite{xrprior1}. W odrębnej pracy oceniali instrukcje AR oparte na początkowym i dynamicznym przydziale tolerancji montażowych \cite{xrprior2}. Wyniki tych badań pokazują, że projekt wskazówek AR i wybrana miara wykonania mają znaczenie dla interpretacji porównań. Nie stanowią jednak bezpośredniego testu przestrzennej instrukcji AR względem ekranowej instrukcji 2D w neutralnych czynnościach domowych.

\subsection{Granica wnioskowania z dostępnych źródeł}
Oba przywołane badania dotyczą precyzyjnych działań manualnych, a ich porównania obejmują odmiany instrukcji AR, nie zaś zaplanowany tu kontrast AR--2D. Nie uzasadniają więc twierdzenia, że w literaturze nie istnieją odpowiednie badania domowe lub że wykazano brak przewagi 2D. Dla oceny luki niezbędne jest odrębne przeszukanie i krytyczne zestawienie prac dotyczących wskazówek AR i ekranowych, z uwzględnieniem typów zadań, definicji pominięcia, sposobów pomiaru i kontekstu uczestników.

\subsection{Zadanie weryfikacyjne przed zatwierdzeniem projektu}
Należy uzupełnić przegląd pierwotnych badań bezpośrednio porównujących AR z instrukcjami 2D w wieloetapowych zadaniach o charakterze przestrzennym. Dopiero po sprawdzeniu metod, wielkości efektów i ograniczeń tych prac będzie można rozstrzygnąć, czy proponowany kontrast uzupełnia rzeczywistą lukę. Na obecnym etapie jest to luka \emph{kandydacka}, nie potwierdzona nowość.

\section{Hipotezy i przewidywania}
\subsection{Kandydat na hipotezę konfirmacyjną}
\paragraph{H1.} W bezpiecznych, wyrównanych zadaniach instrukcje przestrzenne AR zwiększają udział poprawnych kroków względem instrukcji 2D.

H1 odpowiada na PB1 i dotyczy planowanego wyniku głównego: poprawności kroku. Jej kierunek jest przewidywaniem wynikającym z postulowanego ułatwienia lokalizacji wskazówek, a nie ze zweryfikowanej przewagi AR w tym samym typie zadań. Nadanie jej statusu konfirmacyjnego wymaga jeszcze zatwierdzenia operacjonalizacji, estymandu, próby, randomizacji i analizy przed pozyskaniem danych. Do tego czasu H1 pozostaje \emph{kandydatem}, a nie prerejestrowaną hipotezą.

\subsection{Hipotezy eksploracyjne}
\paragraph{H2.} AR zmniejsza pominięcia w krokach o relacjach przestrzennych, ale niekoniecznie przy zadaniach czysto sekwencyjnych.

H2 odpowiada na PB2 i wymaga osobnego zdefiniowania pominięcia oraz uprzedniej klasyfikacji kroków. Sama różnica w jednym typie kroków przy braku istotności w drugim nie dowodzi interakcji. Dla takiego wniosku należałoby bezpośrednio oszacować różnicę efektów między typami kroków.

\paragraph{H3.} Wzrost poprawności nie jest osiągany kosztem przekroczenia wcześniej wybranego limitu czasu ani większej liczby przerwanych sekwencji.

H3 odpowiada na PB3, lecz ma obecnie charakter eksploracyjny, ponieważ limit czasu i dopuszczalny koszt przerwań nie zostały wyznaczone. Brak statystycznie istotnej różnicy czasu nie potwierdza niegorszości. Ewentualne formalne testowanie H3 wymaga uprzedniego określenia marginesów oraz odpowiedniego projektu analizy.

\subsection{Alternatywne rozstrzygnięcia}
Możliwe jest uzyskanie podobnej poprawności w obu warunkach, przewagi instrukcji 2D albo korzystnego efektu AR tylko dla części zadań. Odmienny czas wykonania przy podobnej liczbie błędów również byłby informacyjny. Żaden z tych rezultatów nie został zaobserwowany w ramach tego projektu. Status koncepcji pozostaje \textit{deferred}.

\section{Projekt metodologiczny}
Niniejszy rozdział stanowi \emph{projekt protokołu}, nie dokumentację wykonanego eksperymentu. Nie rozpoczęto rekrutacji, nie zatwierdzono próby, stanowiska ani procedury. Przed badaniem należy zatwierdzić scenariusze, instrukcje, plan analizy i wymagania etyczne oraz utrwalić ich wersje.

\subsection{Schemat badania i jednostka wnioskowania}
Rozważany jest kontrolowany eksperyment porównujący dwa sposoby przedstawienia tej samej instrukcji: ekranowy 2D i przestrzenny AR. Preferowany do oceny wykonalności jest wariant wewnątrzosobniczy, w którym każdy uczestnik realizowałby porównywalne, lecz nieidentyczne sekwencje w obu warunkach. Dobór tego wariantu wymaga pilotażowej oceny transferu uczenia się; alternatywą pozostaje plan międzyosobniczy, jeśli efekt treningu uniemożliwi wiarygodne porównanie. Wybór schematu należy zatwierdzić przed rekrutacją.

Jednostką manipulacji byłby przydział sposobu prezentacji instrukcji do sekwencji zadania, a pojedynczy krok stanowiłby jednostkę zapisu wyniku. Kroki są zagnieżdżone w sekwencjach, a sekwencje mogą powtarzać się u uczestników; analiza musi odzwierciedlać te zależności. Nie wolno utożsamiać liczby kroków z liczebnością niezależnej próby.

\subsection{Populacja docelowa i dobór próby}
Pierwotną populacją odniesienia mogą być dorośli ochotnicy zdolni do zrozumienia instrukcji, samodzielnej obsługi neutralnych przedmiotów i świadomego udziału. Nie zakłada się rekrutacji pacjentów ani osób wymagających opieki. Dobór prawdopodobnie miałby charakter dogodnościowy; zakres uogólniania wyników należałoby ograniczyć odpowiednio do cech rekrutowanej grupy.

Proponowane kryteria włączenia obejmują pełnoletność, świadomą zgodę, zdolność zrozumienia instrukcji i możliwość wykonania zadania w obu porównywanych warunkach, z uwzględnieniem dostępnych udogodnień. Wyłączenia powinny dotyczyć wyłącznie przeciwwskazań bezpośrednio związanych z bezpieczeństwem procedury lub niemożnością uzyskania ważnego pomiaru. Nie należy automatycznie wykluczać osób z niepełnosprawnością ani przyjmować rozpoznań klinicznych bez uzasadnienia. Kryteria przerwania próby i ewentualnego wycofania danych muszą być odrębnie opisane w informacji dla uczestnika.

Liczebność próby badania właściwego nie jest określona. Należałoby ją wyznaczyć na podstawie precyzji oszacowania lub analizy mocy dla przyjętego estymandu, zależności pomiarów, realistycznej zmienności międzyosobniczej i efektu istotnego praktycznie. Wstępna ocena wykonalności nie może być traktowana jako test potwierdzający H1. Rekrutacja i liczebność wymagają zatwierdzenia.

\subsection{Materiał zadaniowy i równoważność warunków}
Zadania powinny wykorzystywać wyłącznie niegroźne przedmioty, takie jak pojemniki, etykiety i elementy do układania, bez zasilania, ostrych krawędzi i gorących powierzchni. Każda sekwencja wymaga jawnej listy kroków, stanu początkowego, oczekiwanego rezultatu, reguły poprawności oraz informacji o tym, czy krok wymaga odniesienia przestrzennego. Zestawy porównawcze należy przygotować przed badaniem, a ich trudność sprawdzić w odrębnej ocenie wykonalności.

W obu warunkach treść semantyczna, kolejność informacji, liczba wskazówek oraz zasady przejścia między krokami powinny być możliwie równoważne. W warunku 2D instrukcja byłaby widoczna na ekranie, a w AR odnosiłaby się przestrzennie do obiektów. Należy z góry opisać sposób wyświetlania błędów, komunikatów pomocy i powtórzeń. Nie wolno poprawiać wyłącznie jednego wariantu na podstawie wyników uczestników badania właściwego.

\subsection{Konstrukty, zmienne i instrumenty}
Zmienną niezależną jest \emph{modalność prezentacji instrukcji}, kodowana jako AR lub 2D. Wynik główny stanowi \emph{poprawność kroku}, oceniana binarnie na podstawie uprzednio przygotowanej rubryki, bez odnoszenia oceny do samej technologii. Poprawny krok oznacza osiągnięcie wymaganego stanu zadania zgodnie z regułą; pominięcie, błędny obiekt i niepoprawna kolejność wymagają osobnych kodów. Trzeba ustalić, czy poprawność po korekcie błędu jest liczona w wyniku głównym, oraz jak traktuje się kroki niewykonane po przerwaniu.

Drugorzędowe obserwacje obejmowałyby liczbę pominięć, czas ukończenia sekwencji, czas kroku, liczbę przerwań, prośby o pomoc i błędy systemowe. Typ kroku (przestrzenny lub sekwencyjny) jest uprzednio definiowaną cechą scenariusza, nie oceną dokonywaną po obejrzeniu wyników. Potencjalnymi moderatorami lub czynnikami zakłócającymi są kolejność warunków, doświadczenie XR, znajomość podobnych zadań, różnice ergonomiczne i jakość zakotwiczenia AR.

Minimalnym instrumentarium byłyby dwie porównywalne implementacje instrukcji, wersjonowane scenariusze, zegar lub rejestr zdarzeń oraz jednolity arkusz oceny kroków. Stanowisko ST-03 jest jedynie propozycją; poprawność działania oprogramowania i sprzętu nie została potwierdzona. EEG, okulometria, VLM i analiza biometryczna nie są wymagane do odpowiedzi na PB1. Ich ewentualne dodanie wymagałoby odrębnego uzasadnienia, zgód i protokołu.

\subsection{Proponowany przebieg sesji}
Po uzyskaniu wymaganych zgód i potwierdzeniu kwalifikacji uczestnik otrzymałby jednolite objaśnienie zadania oraz neutralną próbę zapoznawczą niezaliczaną do analizy. Następnie, zgodnie z uprzednio ustalonym planem przydziału, wykonywałby sekwencje w obu warunkach. Należy zapewnić porównywalne odstępy, odpoczynek i zasady pomocy, a kolejność warunków zrównoważyć lub randomizować zgodnie z zatwierdzonym schematem.

Dla każdej próby należałoby zapisać pseudonimowy identyfikator uczestnika, wersję scenariusza, warunek, kolejność, znaczniki czasu, wyniki kroków i przyczynę przerwania. Po zakończeniu sesji uczestnik powinien mieć możliwość zgłoszenia dyskomfortu oraz wycofania się zgodnie z otrzymaną informacją. Nie przewiduje się adaptacyjnych zmian scenariusza w trakcie zbierania danych potwierdzających.

\subsection{Kontrola jakości i wiarygodność oceny}
Przed pierwszą próbą należy wykonać kontrolę kalibracji przestrzennej, widoczności przedmiotów, czytelności instrukcji, synchronizacji zegarów i kompletności zapisów. Błąd rejestracji wskazówki AR lub utratę logu należy oznaczyć jako zdarzenie techniczne, a nie automatycznie jako błąd uczestnika. Kontrole stanowiska i wersje oprogramowania należy zapisać dla każdej sesji.

Rubrykę oceny poprawności należy sprawdzić na neutralnych przykładach przed badaniem. Przy ocenie obserwacyjnej wskazany jest drugi, niezależny oceniający dla części prób oraz wcześniej określona procedura rozstrzygania niezgodności; możliwość maskowania warunku może być ograniczona przez widoczność interfejsu. Wszelkie rozbieżności, braki danych, wyłączenia prób i odstępstwa od procedury należy zachować w rejestrze, a ich wpływ uwzględnić w analizie wrażliwości. Kryteria jakości i progi akceptacji muszą być określone przed badaniem, bez dopasowywania do uzyskanych wyników.

\subsection{Bezpieczeństwo, etyka i ochrona danych}
Plan musi wykluczać ekspozycję na napięcie elektryczne, wysokie temperatury, ostre przedmioty i konieczność przemieszczania się w niekontrolowanym otoczeniu. Sesję należy przerwać na życzenie uczestnika, przy objawach dyskomfortu lub awarii mogącej ograniczać bezpieczeństwo. Przed rekrutacją konieczna jest właściwa ocena etyczna albo udokumentowane rozstrzygnięcie o jej niewymagalności oraz zatwierdzona informacja i zgoda uczestnika. Tych czynności nie wykonano w ramach niniejszego szkicu.

Zakres zbieranych danych powinien być ograniczony do informacji niezbędnych do oceny zadania; nagrywanie obrazu, głosu i danych biometrycznych nie jest domyślną częścią projektu. Należy ustalić podstawę prawną, okres retencji, dostęp do danych i zasady usunięcia lub wycofania danych. Do repozytorium publicznego nie wolno trafiać z nagraniami ani danymi identyfikującymi uczestników.

\subsection{Otwarte decyzje przed zatwierdzeniem protokołu}
Wymagają rozstrzygnięcia: ostateczny schemat eksperymentu, populacja i liczebność próby, lista scenariuszy i klasyfikacja kroków, szczegółowa rubryka poprawności, reguły czasu i przerwań, realizacja równoważności instrukcji, kryteria braków danych, instrumentarium oraz wymogi etyczne. Żadna z tych propozycji nie jest deklaracją ukończonego pilotażu, gotowości stanowiska ani rozpoczęcia badań.

\section{Plan analizy statystycznej}
\subsection{Estymand i wynik główny}
Kandydackim estymandem jest różnica prawdopodobieństwa poprawnego wykonania kroku pomiędzy AR a 2D, uśredniona w uprzednio określonej populacji zadań i uczestników. Operacyjnie wynik główny można przedstawić jako udział poprawnych kroków przy jawnym mianowniku. Przed analizą należy uzgodnić, czy celem jest kontrast marginalny, oraz przyjąć reguły dotyczące przerwań i braków danych.

\subsection{Model i zależności obserwacji}
Dla binarnej poprawności kroku rozważa się model mieszany z efektem warunku i efektami związanymi z uczestnikiem oraz sekwencją. Kolejność, typ kroku i doświadczenie mogą stanowić wcześniej określone zmienne towarzyszące. Estymatę efektu należy prezentować wraz z przedziałem niepewności; nie można traktować kolejnych kroków tej samej osoby jako niezależnych uczestników.

\subsection{Wyniki uzupełniające i analizy odporności}
Pominięcia, czas wykonania i przerwane sekwencje powinny być raportowane oddzielnie. Różnice według typu kroku można analizować eksploracyjnie. Należy przewidzieć analizę wrażliwości na reguły kodowania nieukończonych sekwencji i brakujących obserwacji. Liczebność, minimalny efekt o znaczeniu praktycznym oraz ewentualny margines niegorszości nie zostały ustalone; nie wolno ich dobierać po poznaniu wyników.

\section{Plan prezentacji wyników}
Jeżeli badanie zostanie zatwierdzone i wykonane, raport powinien rozdzielać liczbę uczestników, sekwencji i kroków od ich wzajemnej zależności. Należałoby przedstawić przepływ prób, kompletność obserwacji, poprawność w obu warunkach, estymowaną różnicę wraz z niepewnością oraz wyniki uzupełniające, w tym przerwania i zdarzenia niepożądane. W obecnym szkicu nie ma wyników, tabel z pomiarami ani wnioskowania statystycznego.

\section{Ramy przyszłej dyskusji}
Ewentualny efekt warunku należałoby interpretować w świetle różnic w kierowaniu uwagą, widoczności instrukcji, uczeniu się i obciążeniu interfejsem. Brak przewagi AR również byłby istotnym rozstrzygnięciem pytania badawczego. Efekt lokalny w neutralnym zadaniu nie uprawnia do twierdzeń o skuteczności klinicznej, samodzielności osób wymagających wsparcia ani zapobieganiu wypadkom domowym.

\section{Ograniczenia koncepcji}
Koncepcja nie ma zatwierdzonego protokołu, potwierdzonej liczebności próby, kompletnej krytycznej syntezy literatury ani udokumentowanej oceny wykonalności stanowiska. Transfer wyników z czynności laboratoryjnych do codziennych zadań i różnych populacji pozostaje otwartą kwestią. Wymaga także kontroli potencjalnego wpływu różnic sprzętowych między AR i 2D.

\section{Wnioski koncepcyjne}
Przedstawiony projekt pozwala sformułować testowalny kontrast między instrukcją przestrzenną AR i ekranową 2D oraz powiązać go z poprawnością kroków. Dalsze prace wymagają weryfikacji literatury, zamknięcia operacjonalizacji, planu doboru uczestników, oceny bezpieczeństwa i decyzji etycznych. Na obecnym etapie nie ma podstaw do wniosków o przewadze AR; status projektu pozostaje \textit{deferred}.

\section*{Informacje autorskie i oświadczenia}
Wskazanym autorem i właścicielem naukowym koncepcji jest Mateusz Pomianek. Afiliacja, identyfikator ORCID i adres do korespondencji wymagają potwierdzenia przed ewentualnym zgłoszeniem manuskryptu.

Wkład autora należy opisać zgodnie z taksonomią CRediT po weryfikacji faktycznie wykonanych czynności. Obecny szkic nie przypisuje nieudokumentowanych ról ani nie przewiduje fikcyjnych współautorów.

Źródła finansowania i ewentualne konflikty interesów nie zostały potwierdzone w dokumentacji tej koncepcji. Odpowiednie deklaracje wymagają ustalenia przed zgłoszeniem; brak takiej informacji nie oznacza deklaracji braku finansowania lub konfliktów.

Nie są dostępne udokumentowane wyniki badań z udziałem ludzi ani dane do publicznego udostępnienia. Przed rozpoczęciem ewentualnego badania konieczne będzie rozstrzygnięcie wymagań dotyczących oceny etycznej, zgody uczestników i przetwarzania danych. Deklarację dostępności danych oraz pozostałe oświadczenia wymagane przez czasopismo należy opracować na podstawie rzeczywistego przebiegu projektu.

% Pola metadanych wydawniczych do potwierdzenia, bez publikowania wartosci domyslnych:
% affiliation_verified:
% orcid_verified:
% corresponding_author_contact_verified:
% CRediT_roles_verified:
% funding_statement_verified:
% conflicts_of_interest_verified:
% ethics_statement_verified:
% data_availability_statement_verified:
% publisher_specific_statements_verified:

\begin{thebibliography}{9}
\bibitem{xrprior1} Zhang, X., He, W., Qin, Y., Billinghurst, M. i in. (2024). Comparison of visual and multisensory augmented reality for precise manual manipulation tasks. \textit{Displays}, 84, 102768. \url{https://doi.org/10.1016/j.displa.2024.102768}.
\bibitem{xrprior2} Zhang, X., He, W., Bai, J., Billinghurst, M. i in. (2025). Evaluation of Augmented Reality instructions based on initial and dynamic assembly tolerance allocation schemes in precise manual assembly. \textit{Advanced Engineering Informatics}, 63, 102954. \url{https://doi.org/10.1016/j.aei.2024.102954}.
\end{thebibliography}
\end{document}

